\documentclass[10pt,a4paper]{amsart}
\usepackage[margin=19mm]{geometry}
\usepackage{amsmath,amssymb,mathtools}
\usepackage[scr=rsfs,cal=euler]{mathalfa}
\usepackage{xfrac}
\usepackage[unicode,hidelinks,bookmarks=false]{hyperref}
\newcommand{\ie}{i.e.\ }
\newcommand{\cf}{cf.\ }
\newcommand{\resp}{respectively\ }
\newcommand{\Subsection}{\subsection}
\providecommand{\nobreakdash}{\nobreak\hbox{-}}
\setlength{\emergencystretch}{2em}
\multlinegap0pt
\usepackage{amssymb}
\usepackage{enumerate}
\usepackage{multirow}
\usepackage{xcolor}
\newtheorem{lemma}{Lemma}
\newcommand{\Fpnmul}{\mathbb{F}_{p^n}^*}
\title{On Differential and Boomerang Properties of a Class of Binomials over Finite Fields of Odd Characteristic}
\author{Namhun Koo, Soonhak Kwon}
\date{Source: arXiv:2506.11486v2 (CC BY 4.0)}
\begin{document}
\maketitle

\noindent\emph{Source and licence.} On Differential and Boomerang Properties of a Class of Binomials over Finite Fields of Odd Characteristic. Version arXiv:2506.11486v2 (CC BY 4.0), distributed by arXiv under the Creative Commons Attribution 4.0 International licence, http://creativecommons.org/licenses/by/4.0/, as recorded in the arXiv licence metadata for this exact version and shown on the abstract page https://arxiv.org/abs/2506.11486v2. Full licence notice: \texttt{licenses/arxiv-2506.11486-CC-BY-4.0.txt}.\par\smallskip

\noindent\textit{Editorial scope.} Counted unit: the article's lemma S10 (source0.tex lines 1156--1163). The article's complete original proof of it is delivered with it (source0.tex lines 1165--1230, one \texttt{proof} environment of 66 lines). No mathematics is written, repaired or re-derived: the statement and the proof are the article's own lines, in the article's own notation, and no retained line is substituted or re-flowed.  Retained uncounted as the source-local context the proof needs: lines 1012--1014; lines 1097--1104.  Nothing was omitted from any retained span, so the package carries no omission record.  No retained line contains a citation, so the wrapper carries no bibliography and no bibliography entry.  No retained line contains a figure environment, an image-inclusion command or drawing code, so no image is delivered and the package declares no asset dependency.  Editorial formatting only, in addition: an AMS \texttt{amsart} preamble that restates the article's own declarations and macros used by the retained lines, namely the environments \texttt{lemma} on separate counters, together with the standard packages those lines need.  The retained environments print in this document as Lemma 1 in order of appearance, whereas the article prints the same items with the numbering of its own full document.  The complete native source of the article as distributed in the arXiv e-print for this exact version is bundled unchanged as \texttt{sources/179795/source0.tex}.
\begin{equation}\label{DU_eqn}
	b=F_{r,u}(x+1)-F_{r,u}(x) = (x+1)^r(1+u\chi(x+1)) - x^r (1+u\chi(x)).
\end{equation}

\begin{lemma}\label{DU_lemma S01}
	Let $u\in \Fpnmul \setminus \{\pm1\}$ and $b\ne 0$. Then, $D_{01}^u(b)\le 2$. Furthermore, $D_{01}^u(b) = 2$ if and only if $\chi(b^2+2(1+u^2))=-1$ and 
	\begin{align*}
		\chi(b(1-u)+(1+u)R_1) = \chi(b(1-u)-(1+u)R_1) &= \chi(1+u^2),\\
		\chi(b(1+u)+(1-u)R_1) = \chi(b(1+u)-(1-u)R_1) &= -\chi(2(1+u^2)),
	\end{align*}
	where $R_1  = (b^2+2(1+u^2))^r$.
\end{lemma}

\begin{lemma}\label{DU_lemma S10}
	Let $u\in \Fpnmul \setminus \{\pm1\}$ and $b\ne 0$. Then, $D_{10}^u(b)\le 2$. Furthermore, $D_{10}^u(b) = 2$ if and only if $\chi(b^2-2(1+u^2))=-1$ and 
	\begin{align*}
		\chi(b(1-u)+(1+u)R_2) = \chi(b(1-u)-(1+u)R_2) &= -\chi(1+u^2),\\
		\chi(b(1+u)+(1-u)R_2) = \chi(b(1+u)-(1-u)R_2) &= \chi(2(1+u^2)),
	\end{align*}
	where $R_2  = (b^2-2(1+u^2))^r $. In particular, if $b^2\in \{(1+u)^2, (1-u)^2\}$, then $D_{10}^u(b)\le 1$. If $D_{10}^u(1+u)= 1$, then $\chi(2(1+u^2))=\chi(u)$. If $D_{10}^u((-1)^{r+1}(1-u))= 1$, then $\chi(2(1+u^2))=-\chi(u)$. 
\end{lemma}

\begin{proof}
If $x\in S_{10}$, then \eqref{DU_eqn} is equivalent to 
\begin{equation}\label{DU_eqn S10}
	b=(1+ u)(x+1)^r- (1-u)x^r .  
\end{equation}
Using the same idea as in Lemma \ref{DU_lemma S01}, we have the following quadratic equation
\begin{equation}\label{DU_eqn S10 quad}
	0=4(1+u^2)^2x^2+4x\left( -2ub^2 +(1+u)^2(1+u^2) \right) +\left(  (1+u)^2 -b^2\right)^2. 
\end{equation}
Thus, we have $D_{10}^u(b) \le 2$. The discriminant of \eqref{DU_eqn S10 quad} is 
\begin{equation*}
	-16b^2(1-u^2)^2\left( b^2-2(1+u^2) \right). 
\end{equation*}
So, \eqref{DU_eqn S10 quad} has two solutions if and only if 
\begin{equation*}
	\chi\left( b^2-2(1+u^2) \right)=-1.
\end{equation*}
In this case, as in Lemma \ref{DU_lemma S01}, two solutions are
\begin{equation*}
	x_{10}= -\frac{\left( b(1-u)-\epsilon (1+u)R_2\right)^2}{4(1+u^2)^2},
\end{equation*}
where $\epsilon\in \{1,-1\}$, and we have
\begin{equation*}
	x_{10}+1 =\frac{\left( b(1+u)+\epsilon(1-u)R_2\right) ^2}{4(1+u^2)^2}.
\end{equation*}
We can see that $\chi(x_{10})=-1$ and $\chi(x_{10}+1)=1$. Substituting $x=x_{10}$ in the right hand side of \eqref{DU_eqn S10},
\begin{align*}
	(1+ u)&(x_{10}+1)^r- (1-u)x_{10}^r \\
	&=\frac{b}{2(1+u^2)}\left(  (1+u)^2\chi\left( \frac{ b(1+u)+\epsilon(1-u)R_2}{2(1+u^2)}\right) -(-1)^r(1-u)^2\chi\left( \frac{ b(1-u)-\epsilon(1+u)R_2}{2(1+u^2)}\right) \right) \\
	&+\frac{(1-u^2)R_2 \epsilon}{2(1+u^2)}\left(  \chi\left( \frac{ b(1+u)+\epsilon(1-u)R_2}{2(1+u^2)}\right) +(-1)^r\chi\left( \frac{ b(1-u)-\epsilon(1+u)R_2}{2(1+u^2)}\right) \right).
\end{align*}
Let us denote the resulting expression by $E_{10}$. Then,
\begin{align*}
	\chi\left( \frac{ b(1+u)+\epsilon(1-u)R_2}{2(1+u^2)}\right)=1, \ \chi\left( \frac{ b(1-u)-\epsilon(1+u)R_2}{2(1+u^2)}\right) = (-1)^{r+1}\ &\Rightarrow\ E_{10}=b\\
	\chi\left( \frac{ b(1+u)+\epsilon(1-u)R_2}{2(1+u^2)}\right)=-1, \ \chi\left( \frac{ b(1-u)-\epsilon(1+u)R_2}{2(1+u^2)}\right) = (-1)^r\  &\Rightarrow\ E_{10}=-b\\
	\chi\left( \frac{ b(1+u)+\epsilon(1-u)R_2}{2(1+u^2)}\right)=1,\ \chi\left( \frac{ b(1-u)-\epsilon(1+u)R_2}{2(1+u^2)}\right) = (-1)^r\ &\Rightarrow\ E_{10}=\frac{4ub+2\epsilon (1-u^2)R_2}{2(1+u^2)}\\
	\chi\left( \frac{ b(1+u)+\epsilon (1-u)R_2}{2(1+u^2)}\right)=-1,\ \chi\left( \frac{ b(1-u)-\epsilon (1+u)R_2}{2(1+u^2)}\right) = (-1)^{r+1},\ &\Rightarrow\ E_{10}=\frac{-4ub- 2\epsilon (1-u^2)R_2}{2(1+u^2)}
\end{align*}
In the third and fourth cases above, we obtain
\[
E_{10}=\pm \frac{2ub+\epsilon(1-u^2)R_2}{1+u^2}.
\]
A direct computation shows that $E_{10}=b$ holds if and only if
$b^2=(1+u)^2$ in the third case, and $b^2=(1-u)^2$ in the fourth case.


If $b^2 = (1+u)^2$, then $R_2=(-(1-u)^2)^r=(-1)^r (1-u)\chi(1-u)$. Suppose that $b=1+u$. If $\epsilon = (-1)^{r+1} \chi(1-u)$, then
\begin{equation*}
	x_{10} = -\frac{\left( (1+u)(1-u)-\epsilon (-1)^r(1+u)(1-u)\chi(1-u)\right)^2}{4(1+u^2)^2} = -\frac{(1-u^2)^2}{(1+u^2)^2}.
\end{equation*}
If $\epsilon = (-1)^{r} \chi(1-u)$, then we have $x_{10}=0\not \in S_{01}$. Thus, we can see that $D_{01}^u(1+u)\le 1$. Moreover, if $D_{01}^u(1+u)= 1$, then
\begin{equation*}
	1=\chi\left(\frac{ b(1+u)+\epsilon(1-u)R_2}{2(1+u^2)}\right)=\chi\left( \frac{ (1+u)^2-(1-u)^2}{2(1+u^2)}\right) = \chi\left( \frac{4u}{2(1+u^2)}\right), 
\end{equation*}
or equivalently, $\chi(2(1+u^2))=\chi(u)$. Similarly, we have $D_{01}^u(-(1+u))\le 1$.

If $b^2=(1-u)^2$, then $R_2=(-(1+u)^2)^r=(-1)^r (1+u)\chi(1+u)$. Suppose that $b=(-1)^{r+1}(1-u)$. If $\epsilon = -\chi(1+u)$, then
\begin{equation*}
	x_{10}= -\frac{\left( (-1)^{r+1}(1-u)^2-\epsilon (-1)^r(1+u)^2 \chi(1+u) \right)^2}{4(1+u^2)^2} = -\frac{4u^2}{(1+u^2)^2}.
\end{equation*}
If $\epsilon = \chi(1+u)$, then we have $x_{10}=-1\not \in S_{01}$. Thus, we can see that $D_{01}^u((-1)^{r+1}(1-u)) \le 1$. Furthermore, if $D_{01}^u((-1)^{r+1}(1-u)) = 1$,  then 
\begin{equation*}
	(-1)^{r+1}=\chi\left( \frac{ b(1-u)-\epsilon(1+u)R_2}{2(1+u^2)}\right) =\chi\left( \frac{ (-1)^{r+1}(1-u)^2-(-1)^{r+1}(1+u)^2}{2(1+u^2)}\right)  =(-1)^{r}\chi\left( \frac{4u}{2(1+u^2)}\right), 
\end{equation*}
or equivalently, $\chi(2(1+u^2))=-\chi(u)$. Similarly, we have $D_{01}^u((-1)^{r}(1-u)) \le 1$.
\end{proof}

\end{document}
